\documentclass{article}
\usepackage{amsmath, amssymb, amsthm}
\usepackage{hyperref}
\usepackage{geometry}
\geometry{margin=1in}

\title{Erd\H{o}s Problem \#406}
\date{Last updated: 2025-08-31}
\author{Source: erdosproblems.com}

\begin{document}
\maketitle

\noindent\textbf{Status:} OPEN \quad \textbf{No prize}

\noindent\textbf{Tags:} number theory, base representations

\medskip

\noindent\textbf{Problem Statement:}

Is it true that there are only finitely many powers of $2$ which have only the digits $0$ and $1$ when written in base $3$?




The only examples seem to be $1$, $4=1+3$, and $256=1+3+3^2+3^5$. If we only allow the digits $1$ and $2$ then $2^{15}$ seems to be the largest such power of $2$. This would imply via Kummer's theorem that\[3\mid \binom{2^{k+1}}{2^k}\]for all large $k$. Saye \cite{Sa22} has computed that $2^n$ contains every possible ternary digit for $16\leq n \leq 5.9\times 10^{21}$. Let $N(x)$ count the number of $n\leq x$ such that $2^n$ has only the digits $0$ and $1$ in base $3$. Narkiewicz \cite{Na80} proved\[N(x)\leq 1.62 x^{\log_32}\]This is mentioned in problem B33 of Guy's collection \cite{Gu04}. There are many generalisations possible - see, for example, \cite{AbLa14} and \cite{La09}.




References

[AbLa14] Abram, William C. and Lagarias, Jeffrey C., Intersections of multiplicative translates of 3-adic {C}antor sets . J. Fractal Geom. (2014), 349--390.

[Gu04] Guy, Richard K., Unsolved problems in number theory . (2004), xviii+437.

[La09] Lagarias, Jeffrey C., Ternary expansions of powers of 2 . J. Lond. Math. Soc. (2) (2009), 562--588.

[Na80] Narkiewicz, W., A note on a paper of {H}. {G}upta concerning powers of two and three: ``{P}owers of {$2$}\ and sums of distinct powers of {$3$}''\ [{U}niv. {B}eograd. {P}ubl. {E}lektrotehn. {F}ak. Ser. {M}at. {F}iz. {N}o. 602-633 (1978), 151--158 (1979);\ {MR} 81g:10016] . Univ. Beograd. Publ. Elektrotehn. Fak. Ser. Mat. Fiz. (1980), 173--174.

[Sa22] Saye, Robert I., On two conjectures concerning the ternary digits of powers of two . J. Integer Seq. (2022), Art. 22.3.4, 9.


\bigskip
\noindent\small{Source: \url{https://www.erdosproblems.com/406}}
\end{document}
